\chapter*{Introduction générale}
\addcontentsline{toc}{chapter}{Introduction générale}
\markboth{Introduction générale}{Introduction générale}

L'industrie automobile mondiale traverse actuellement une phase de transformation numérique et technologique sans précédent. Dans un contexte fortement concurrentiel, marqué par l'émergence des véhicules électriques, des systèmes de conduite autonome et des normes de sécurité de plus en plus drastiques, les équipementiers automobiles doivent adapter leurs outils de production et de gestion pour rester performants. Le Groupe COFAT, acteur historique et incontournable de ce secteur depuis sa création en 1984, n'échappe pas à cette dynamique. Avec ses 5 600 collaborateurs répartis sur 10 entités à travers 6 pays (Tunisie, Allemagne, Égypte, Mexique, Brésil et États-Unis), COFAT s'est imposé comme un leader dans la conception, le développement et la fabrication de faisceaux de câbles automobiles. Ses clients, qui incluent des géants tels que Stellantis, Volkswagen, Renault, Scania, Volvo et Daimler, exigent une qualité irréprochable et une réactivité immédiate face aux variations de la demande.

La gestion de la capacité de production à une échelle aussi vaste représente un défi majeur. Chaque site de production (PLANT) doit continuellement s'assurer qu'il dispose des équipements adéquats, de l'espace physique nécessaire et des ressources humaines suffisantes pour honorer les commandes. Jusqu'à présent, la planification de ces capacités et la consolidation des besoins en investissement étaient réalisées de manière semi-manuelle, reposant massivement sur des échanges de fichiers Excel entre les différents sites et le siège. Ce processus chronophage, sujet aux erreurs de saisie et manquant de traçabilité, ralentissait considérablement la prise de décision, en particulier lors de la validation des budgets d'investissement non industriels et industriels. Face à ce constat, le besoin d'un système centralisé, automatisé et collaboratif est devenu impératif. 

C'est dans ce cadre précis que s'inscrit mon Projet de Fin d'Études (PFE), intitulé \textbf{\textit{Cofat Capacity Study}}. Il s'agit d'un système global de gestion et de consolidation des capacités du groupe COFAT. Cette plateforme web sur mesure, développée avec des technologies modernes (React.js en frontend, Node.js/Express en backend et Microsoft SQL Server en base de données), vise à digitaliser intégralement le processus de planification et de budgétisation. En offrant un portail multi-utilisateurs sécurisé, l'application permet aux gestionnaires de sites de soumettre leurs données locales via des imports Excel automatisés, tandis que le moteur de consolidation agrège ces informations en temps réel à l'échelle mondiale. De plus, un workflow budgétaire transparent facilite la collaboration entre les demandeurs et le département des achats.

Le présent rapport détaille l'ensemble des phases de conception et de réalisation de ce projet d'envergure, en suivant une méthodologie agile (Scrum). Il est structuré en quatre chapitres principaux, organisés comme suit :

\begin{itemize}
    \item \textbf{Le Chapitre 1 : Cadre du projet.} Ce chapitre dresse la présentation de l'organisme d'accueil COFAT ainsi que le cadre général du projet. Il détaille l'étude de l'existant à travers une analyse comparative illustrée de trois solutions du marché et de la solution interne historique, puis expose la synthèse du choix de la solution et la méthodologie de projet Scrum/UML adoptée.
    \item \textbf{Le Chapitre 2 : Analyse et spécification des besoins.} Ce chapitre est consacré à l'identification des trois rôles réels (Site Manager, Admin, Achat), à la spécification des besoins fonctionnels et non fonctionnels, à la modélisation UML globale (Cas d'utilisation, Diagramme de classes), ainsi qu'à la planification agile du travail (deux Releases et quatre Sprints du Product Backlog), à l'architecture 3-tiers et à l'environnement de travail.
    \item \textbf{Le Chapitre 3 : Release 1 : Socle technique, authentification et planification locale.} Ce chapitre présente la réalisation de la Release 1, couvrant le Sprint 1 (architecture 3-tiers, authentification sécurisée JWT et administration du référentiel des sites) et le Sprint 2 (planification locale des Équipements, des Espaces et des RH avec moteur d'importation Excel).
    \item \textbf{Le Chapitre 4 : Release 2 : Consolidation groupe, catalogue standard et workflow budgétaire.} Ce chapitre détaille la réalisation de la Release 2, couvrant le Sprint 3 (moteur de consolidation mondiale Cofat Group et catalogue centralisé StandardInvestment) et le Sprint 4 (workflow collaboratif du budget non-industriel, restriction du rôle Achat à 2 colonnes, système de notifications temps réel et exportation de rapports Excel/PDF).
\end{itemize}
